\section{U003 Mixed precision layout and scales}
\label{sec:u003}

A matrix accelerator's ability to multiply a low-bit-width operand does not by itself explain how that operand is produced, represented, or delivered. An attention kernel may form scores in FP32, keep a softmax intermediate in BF16, and feed a subsequent product with packed FP4 values and separately encoded scales. A weight-only kernel may store INT4 weights while presenting a floating-point operand to its matrix engine. These paths have different numerical meanings even when the stored payloads occupy the same number of bits. The relevant question is therefore: after competent compilation, which numerical conversions, scale computations, layout changes, and producer--consumer ownership changes remain, and who performs them?

This question connects storage lifetime in Section~\ref{sec:u002} to reduction execution in Section~\ref{sec:u005}. Here the concern is the representation presented to a reduction or matrix operation; the scheduling and throughput of its internal arithmetic are separate. The comparison includes NVIDIA GPU, Ascend, TPU, Tenstorrent, x86 AVX, Arm NEON/SVE, and the proposed superscalar NPU (SSNPU). Product generations, public APIs, virtual instructions, and implementation observations are identified explicitly. None of the examples is a matched performance measurement.

\subsection{Six contracts behind a low-precision tensor}

An implementation must preserve six related contracts. The first is the logical tensor: indices, valid extents, and reduction axes. The second is physical storage: strides, blocking, padding, and swizzles. The third specifies the bit positions of subbyte values. The fourth assigns values to threads, registers, banks, or local-memory cells. The fifth defines numerical input, intermediate, accumulator, and output types, including rounding, saturation, and exceptional values. The sixth maps each value to its scale and, where applicable, zero point. A layout name only describes part of this information. A storage reinterpretation can change a view without performing the quantization expressed by an arithmetic conversion. MLIR's quantization dialect makes the distinction between an expressed value and its storage representation explicit~\citep{u003_mlir_quant}.

For an illustrative symmetric integer quantizer, let $G_g$ be a group of source elements and choose a positive scale $s_g$. With a signed $b$-bit storage range $[q_{\min},q_{\max}]$, the operation is
\begin{equation}
 q_i=\operatorname{clip}_{[q_{\min},q_{\max}]}
       \left(\operatorname{round}_{\rho}(x_i/s_g)\right),
 \qquad \widehat{x}_i=s_gq_i,\qquad i\in G_g.
 \label{eq:u003_quant}
\end{equation}
The rounding rule $\rho$, scale selection, and treatment of an all-zero group are part of the algorithm. An asymmetric quantizer adds a zero point, and a floating-point quantizer replaces the integer rounding operation with rounding to a specified floating-point format. Neither is obtained merely by changing the destination's bit-width annotation. In particular, integer INT4, E2M1 FP4, and a block-floating representation with shared exponents are not interchangeable.

The scale's position relative to a reduction also matters. Per-row scales for $A$ and per-column scales for $B$ permit
\begin{equation}
 C_{ij}\simeq s^A_i s^B_j\sum_k q^A_{ik}q^B_{kj}.
 \label{eq:u003_outer_scales}
\end{equation}
If scales change along the reduction dimension, with blocks $K_g$, the corresponding expression becomes
\begin{equation}
 C_{ij}\simeq\sum_g s^A_{ig}s^B_{gj}
                     \sum_{k\in K_g}q^A_{ik}q^B_{kj}.
 \label{eq:u003_inner_scales}
\end{equation}
There is generally no single final multiplier that can replace these block-specific products. Native scaled matrix instructions can place this work inside the matrix execution path, subject to their legal formats and scale layouts. They do not make scale production, transport, or representation disappear.

\subsection{A worked preparation path and its storage accounting}

Consider an illustrative $32\times64$ output fragment with FP32 accumulated values and a packed four-bit destination. Ignoring alignment and padding, the accumulator payload is $32\cdot64\cdot4=8192$ bytes, while the packed payload is $32\cdot64/2=1024$ bytes. A one-byte scale for every row and every 64-column group contributes 32 more bytes; a one-byte scale per 32-column group contributes 64 bytes. These are payload counts, not a claim that either scale policy achieves a particular accuracy. A BF16 intermediate for the same fragment occupies another 4096 bytes if it is materialized. It may overlap other lifetimes or be eliminated by fusion, but the comparison must establish which occurs.

The pinned SuperNPUBench HiF4 level-one quantization example provides a concrete preparation sequence. It starts from FP32 matrix results, uses BF16 intermediates, derives and encodes E6M2 scales, applies broadcast scaling, assembles two $32\times32$ fragments into a $32\times64$ Tile, and converts to packed HiF4X2 output. The two logical values per byte affect output address calculations. The scale carrier also has its own physical descriptor geometry: one meaningful scale per row does not imply that its allocated carrier has exactly one physical column. This is a source observation for level-one quantization, rather than evidence for an entire multilevel HiF4 algorithm~\citep{u003_hif4_bench}.

Three operations in this sequence should not be collapsed into the word ``shuffle.'' Assembly joins logical fragments and coordinates their lifetimes. Conversion changes numerical representation and packs the result. A transpose or layout-aware transfer changes where values reside. The last two may sometimes be fused, and assembly can sometimes be avoided by choosing a different producer decomposition, but removing one does not prove that the others are free. In the separate low-precision attention example, probabilities are converted into E2M1x2 fragment slots with E8M0 scales for 32-column groups before assembly for the $PV$ product. That scale contract is different from the HiF4 example's E6M2 scale per 64 columns~\citep{u003_lowp_attention}.

This distinction also prevents a misleading bandwidth argument. Reducing the final payload from 8192 to 1024 bytes does not guarantee an eightfold reduction in total traffic. Scale traffic, index or descriptor traffic, intermediate reads and writes, padding, and format-specific movement may all contribute. Conversely, charging every conceptual intermediate as an external-memory round trip is equally unfair: a compiler can keep intermediates local, fuse producers and consumers, or remove their materialization altogether.

\subsection{The baseline is an optimized representation pipeline}

A meaningful baseline first propagates compatible layouts through producers and consumers. If two representations differ only in a view or in thread-local register naming, conversion need not move data. If a value must move between distinct owners, however, the implementation needs an exchange mechanism. On a GPU this can be a register shuffle or shared-memory exchange; in other machines it can be a vector permutation, local-memory movement, or a rearrangement in an engine's input path. Triton's layout model is useful precisely because it records ownership rather than assuming that a tensor's logical dimensions determine its physical distribution~\citep{u003_triton_layouts}.

Fusion and amortization are separate optimizations. A fused matrix epilogue can apply scaling, clamping, and conversion while accumulator data are available. CUTLASS documents the relationship between accumulator distribution, epilogue exchange, and efficient output stores~\citep{u003_cutlass_gemm}. Static weights may also be prepacked into an instruction-friendly order and reused for many invocations. Prepacking is advantageous only after its setup and additional stored representation are charged to the reuse interval. A transient activation with a new dynamic scale each invocation offers a different opportunity.

Vectorization and unrolling introduce another tradeoff. A compiler can combine conversions with packing, select a consumer-compatible arrangement, and use unrolling to expose independent work or reuse a shared load. It must also account for live-register growth, spills, code size, and the profitability of mixed-type operations. LLVM's vectorizers explicitly distinguish legal transformations from profitable ones~\citep{u003_llvm_vectorizers}. A semantic Tile interface can reduce decomposition performed solely to match a fixed physical width; it does not remove reuse-driven unrolling, algorithmic blocking, or the need to choose numerical representations.

The accounting vocabulary should consequently be precise. A removed redundant conversion is \emph{eliminated}. A static packing pass reused over many executions is \emph{amortized}. A conversion overlapped with independent matrix work is \emph{hidden}, but still occupies arithmetic, ports, and storage. A necessary scale reduction or ownership exchange is \emph{remaining}. Combining these outcomes into one instruction-count reduction obscures where the actual resource savings occur.

\subsection{NVIDIA: packed conversion, fragment ownership, and native scaling}

PTX exposes packed floating-point conversions, but is a virtual ISA rather than final SASS. A supported FP32-to-E4M3x2 conversion produces two FP8 values in a 16-bit carrier; supported FP32-to-E2M1x2 forms similarly pack two FP4 values. The rounding and finite-saturation modifiers define numerical behavior, while the input must already carry the intended scaling. Packing two values does not establish a complete matrix-operand layout transformation. Hopper WGMMA offers FP8 matrix paths with prescribed operand layouts and FP32 accumulation; its sign controls must not be mistaken for arbitrary quantization scales~\citep{u003_ptx}.

Blackwell adds native block-scaled paths with explicit restrictions. The SM100-oriented CUTLASS documentation distinguishes MX formats using UE8M0 scales for 32-element groups from NVFP4 using UE4M3 scales for 16-element groups. These scale values require the prescribed storage organization. SM120 paths are separately described, so the SM100 \texttt{tcgen05} interface is not a generic instruction description for all Blackwell products~\citep{u003_cutlass_blackwell}. Libraries can implement scale policies on hardware without an identical native operand form; the proper comparison is the resulting optimized computation, not the presence or absence of a particular mnemonic.

Ownership remains the central cost question. Matrix instructions determine a fragment distribution; a following row reduction or output store may prefer another distribution. Software and the compiler choose tiling, warp roles, shared-memory layouts, and synchronization. Hardware supplies instruction execution, warp readiness, register and memory access, and asynchronous completion mechanisms. A compatible producer--consumer layout can avoid an exchange, while an incompatible one may require movement and additional live registers. The potential advantage of native scaling is largest when its group size, scale representation, and operand layout already match the algorithm. If they do not, the conversion into that native representation belongs in the cost of using it.

\subsection{Ascend and TPU: legal conversion paths and compiler-selected precision}

Ascend's public interfaces expose distinct conversion, transfer, matrix, and postprocessing responsibilities. The CANN 9.1 \texttt{Cast} tables are generation-specific: A2/A3 include half-to-INT4 conversion, but the table is not evidence of a direct FP32-to-INT4 operation. Newer 950 tables add other combinations, including BF16-to-packed-FP4. A staged conversion must preserve the intended rounding behavior; FP32-to-BF16-to-FP4 can differ from a single rounding to the final format. Cast mask and repetition parameters depend on the larger element width, and subbyte operands impose paired-mask and count constraints~\citep{u003_ascend_cast}.

On the 950 family, \texttt{MmadMx} supplies a different contract: supported FP4 pairs or supported FP8 pairs, E8M0 scale data, and FP32 result storage. Scales are grouped along $K$ in sets of 32. The documented matrices and scale buffers have specific physical arrangements and alignment constraints; independently supported FP4 and FP8 families do not imply every mixed-family pair~\citep{u003_ascend_mmadmx}. The kernel or compiler still chooses its decomposition and stages data through the relevant UB/L1/L0 resources. Hardware executes the vector, transfer, Cube, and postprocessing operations. Format legality, local capacity, events, and scale loads bound the available overlap. API documentation is sufficient to establish these responsibilities without pretending to reveal an undocumented raw encoding. The public 3510 architecture preview additionally describes SIMD register-resident intermediates and SIMT hardware with warp schedulers. It also describes Fixpipe channel merging/splitting and supported layout conversion. These generation-specific mechanisms reinforce that staging and packing should be assessed on the selected path, rather than assuming every intermediate returns to UB~\citep{u003_ascend_3510}.

TPU's Pallas matrix-multiplication tutorial illustrates BF16 input data, an FP32 accumulator reference, and conversion to BF16 output~\citep{u003_pallas_matmul}. These are three different choices. Requesting an FP32 result or accumulator does not alone establish multiplication of unrestricted FP32 operands at full FP32 precision. Pallas precision controls govern how supported operations are lowered, with generation-dependent arithmetic behavior~\citep{u003_pallas_details}. Public compiler IR reveals matrix and conversion intent, not final machine instructions.

Current experimental interfaces broaden the set of explicitly controlled matrix paths: the recorded JAX sources include FP8 cases on newer TPU generations and INT8/INT4 FIFO-oriented cases on earlier supported generations. These are scoped compiler/API observations, not evidence that every installed kernel supports every combination~\citep{u003_jax_explicit}. The general ownership split is that the compiler chooses supported precision, layout, and scheduling; the runtime supplies buffers and launch context; and hardware executes matrix, vector, and transfer work. Scale materialization, VMEM intermediates, padded blocks, and conversion remain real costs. A compile-time array shape is compatible with runtime values and addresses, an important distinction revisited in Section~\ref{sec:u004}.

\subsection{Tenstorrent: storage formats do not directly name matrix arithmetic}

Tenstorrent's unpack--math--pack organization makes representation boundaries particularly visible. Data in L1 are unpacked into source registers, matrix or vector work writes destination registers, and packing returns results to L1. Input format configuration and output packing configuration are distinct controls. Destination accumulation width, including a supported FP32 destination mode, does not by itself specify full-FP32 multiplication. Circular-buffer ownership and engine synchronization surround these operations~\citep{u003_tt_engines}.

For the documented Wormhole path, BFP8 and BFP4 are block-floating storage representations with shared exponents over groups of 16 values. Unpacking interprets those values for the supported matrix input paths. BFP4 must therefore not be substituted for E2M1 FP4 or ordinary INT4 simply because all use a four-bit label somewhere in their representation~\citep{u003_tt_unpack}. Rounding can occur at both input and output boundaries, and fidelity phases are additional execution work. The documented matrix instruction's address and source controls rely on accompanying format and engine state~\citep{u003_tt_matrix}.

The algorithm chooses whether the format's error and exponent-sharing behavior are acceptable. The compiler and kernel choose tiles, format transitions, buffer lifetimes, and the unpack/math/pack schedule. The hardware provides those engines, destination storage, and movement control. The benefit is strongest when data arrive in a reusable supported format and adjacent operations retain compatible engine state. Frequent transitions, poorly filled tiles, or long-lived destination values can instead consume control bandwidth and capacity. This account does not infer support or absence of unrelated low-precision formats on every Tenstorrent generation.

\subsection{x86 and Arm: fixed width, scalable width, and narrowing placement}

Fixed-width vector instructions make element count and destination occupancy visibly different. On an appropriate x86 AVX-512 target, \texttt{VCVTPS2PH} converts 16 FP32 values in a ZMM source into 16 FP16 values in a YMM result. The operation preserves 16 logical elements rather than filling a 512-bit narrow result. Other supported sequences can convert FP16 to integer words and then saturating-narrow to bytes, but that is integer conversion rather than FP8 quantization; scaling and range policy are still required~\citep{u003_intel_sdm}.

Packing order can matter as much as width. For AVX2 \texttt{VPACKSSWB}, two 16-word inputs $A$ and $B$ produce bytes in 128-bit-lane groups:
\begin{equation}
 [A_0\ldots A_7,\;B_0\ldots B_7,\;A_8\ldots A_{15},\;B_8\ldots B_{15}].
 \label{eq:u003_avx_pack}
\end{equation}
A consumer expecting all of $A$ followed by all of $B$ requires rearrangement, whereas a compatible consumer can retain this order. The existence of a reduction intrinsic likewise does not prove that one machine opcode performs it. The compiler's opportunity is to propagate useful lane arrangements across the entire computation rather than repeatedly normalize them~\citep{u003_intel_permute}.

Arm NEON supplies fixed 128-bit vectors. Its narrowing operations distinguish the lower and upper portions of a destination; paired narrowing instructions can combine two sources. The numerical conversion and scale formula remain separate~\citep{u003_arm_neon}. SVE/SVE2 instead support vector-length-agnostic programming with predicates and implementation-dependent vector length. This changes loop portability and tail handling, but does not guarantee that each conversion fills every destination position.

SVE floating-point narrowing can preserve the source element positions in wider slots, while SVE2 bottom/top integer narrowing places results in alternating narrow positions. A consumer may use that arrangement directly or require a permutation. Narrowing stores provide another route: once numerical range and representation are correct, storing the low byte from each wider element can avoid materializing a fully packed vector register~\citep{u003_arm_sve_examples,u003_arm_acle}. These are concrete alternatives to assuming that every width change requires a fixed packing sequence. They also explain why vector-length agnosticism cannot be equated with free conversion or guaranteed useful-lane occupancy. Ordinary SVE's one-dimensional vectors and reductions are distinct from a two-dimensional Tile-axis contract; SME's ZA is a separate architectural interface.

\subsection{SSNPU: typed geometry shifts decomposition into bounded microsteps}

The SSNPU proposal exposes supported typed, shaped Tiles and lets the backend sequence physical work. A Tile's logical extent can change without changing the instruction encoding length. Physical capacity, valid rows and columns, and layout legality remain distinct: runtime valid extent does not create unlimited runtime storage. The inspected Linx API carries dtype, layout, geometry, and discrete capacity choices, while its conversion emitter encodes source type and destination representation/rounding fields~\citep{u003_linx_geometry,u003_linx_emit}.

The pinned API's CUBE cell is 128 bytes. For a cell with $R$ rows and $b$ bits per element, its column capacity is
\begin{equation}
 C_{\mathrm{cell}}=\frac{128\cdot8}{Rb}.
 \label{eq:u003_cell_geometry}
\end{equation}
Thus M32 cells hold one, two, four, or eight columns at 32, 16, 8, or 4 bits respectively; M16 cells hold twice those column counts. This is physical geometry arithmetic, not a throughput estimate. Logical columns must be padded to legal cell boundaries, with capacity rounding in addition. The four-bit M16 mapping also includes a local column-group arrangement, so a naive row-major nibble formula is not a universal address rule~\citep{u003_linx_geometry}.

Typed row and column operations can communicate an axis without making software enumerate every physical lane. However, the inspected row/column sum interfaces do not expose an independent accumulator-width argument. An FP32 matrix accumulator elsewhere in the program cannot be used to infer widened accumulation for an FP16 reduction. Typed axis APIs also exist in other systems; the design opportunity is their integration with the chosen storage, dependency, and issue model rather than ownership of the concept itself~\citep{u003_linx_emit}.

At the recorded TimingSim revision, vector operations are split into cell-sized steps. Widening can map one source cell to multiple destination writes, while narrowing targets packed destination cells across source steps. The review also records a remaining whole-cell-overwrite issue in the narrowing value model and absence of numerical \texttt{TCVT} execution in that inspected compute path. Consequently, sequencing and footprint evidence cannot establish packed-value correctness. These are implementation observations at revision \texttt{d5419b3}, not restrictions inferred for the entire architecture and not RTL or silicon validation.

The proposed benefit is a reduction in explicit width decomposition and some producer--consumer coordination for supported operations. The algorithm still owns numerical policy and scale groups; the compiler still owns legal layouts, blocking, and profitable fusion; the runtime still provides the execution context and memory resources. Hardware or its model must manage cell sequencing, tags and generations, readiness, capacity admission, destination assembly, bank access, and release. Larger semantic Tiles reduce neither the number of values that must be converted nor the amount of live payload by definition. They can help when a compatible multi-stage pipeline would otherwise expose substantial width and ownership bookkeeping; they can be neutral or costly when extra control, padding, or layout transitions dominate.

\subsection{What a convincing comparison must establish}

A comparison should begin with one numerical task and its acceptable error, then state storage, matrix operand, accumulator, and final output types independently. It should specify scale groups and scale encoding, conversion behavior, exceptional values, physical layout, and valid tails. The optimized path should identify which intermediates disappear, which formats are reused, which operations overlap, and which remain on the critical path.

The resulting ledger includes bytes and lifetimes, scale and conversion work, movement between owners, register or local-buffer pressure, padding, and control state. A lower visible instruction count is useful only in relation to those quantities. No source considered here supplies a matched quantitative advantage for the proposed SSNPU, and no evidence supports a ranking of compiler effort, utilization, or PPA across these architectures. The defensible claim is narrower and more useful: typed geometry and native format support can relocate preparation work and coordination, while numerical meaning and physical representation continue to constrain every implementation.
